% 本文件由 LaTeX 排版 Agent 根据有效 translation.json 自动生成。
% author=Methodius; work_id=0628; title=论圣咏

\chapter{论圣咏}

% structure: p0001 source=h1 target=omitted reason=page-title-already-rendered-by-parent

一、愿天主受赞颂！让我们前行，弟兄们，从奇迹到主的奇事，仿佛从力量到力量。因为正如一条金链，其环节紧密相连、彼此勾锁，一环牵动另一环，并与之结合，从而延续金链——同样，由神圣福音先后传递的奇迹，也引导并滋润天主的教会——这教会以欢庆为乐——所用的食物不是那会朽坏的，而是那存留到永生的。\BibleRef{若 6:27} 那么，来吧，诸位蒙爱者，让我们也以准备好的心、侧耳倾听，听主、我们的天主透过先知和福音论及这至圣庆节将对我说什么。的确，祂必将平安启示给祂的子民、祂的圣徒，以及那些将心转向祂的人。今天，先知的号角声唤醒了世界，使散居万邦的天主教会充满喜乐与欢愉。并且，他们将信友从神圣斋戒的操练中，以及从那对抗肉身私欲的\OriginalTerm{训练场}{palaestra}上召来，教导他们向赐予胜利的基督唱一首得胜的新赞美诗与和平的新歌。那么，来吧，每一个人，让我们在主内喜乐；哦，来吧，万民，让我们拍掌，以欢歌向天主我们的救主呼喊。愿无人在这恩宠中无份；愿无人错过这召唤；因为悖逆之子的命运已定，归于灭亡。——愿无人疏忽迎接君王，免得被关在新郎的洞房外。——愿我们中无人以愁容接待祂，免得与那些邪恶的市民一同受罚——我指的是拒绝接受主为他们君王的市民。\BibleRef{路 19:27} 让我们全体欢然同聚；让我们欣然接待祂，以最大的虔诚庆祝我们的节庆。不用衣袍，却将我们的心铺在祂面前；以圣咏和赞美诗，向祂献上感恩的欢呼；并且不住地高呼：\Quote{奉主名而来的，当受赞颂；} 因为祝福祂的人必蒙祝福，咒骂祂的人必受诅咒。\BibleRef{创 27:29} 我还要再说，也不会停止劝勉你们行善：来吧，蒙爱者，让我们赞美那位当受赞颂者，好使我们自己也蒙祂赞颂。这篇讲道召唤各个年龄和身份的人都来赞美主：地上的君王和万民；首领和地上的一切判官；少年和少女——而在这奇迹中，新意在于：婴孩和吃奶者的柔弱天真之龄，获得了用感恩的宣认向天主高唱赞美诗的首要位置，这赞美诗正是天主教导他们的，其曲调是梅瑟从前带领百姓出埃及时所唱的——即：\Quote{奉主名而来的，当受赞颂。}\par

二、今天，圣达味极其喜乐，被孩童夺去了他的\OriginalTerm{竖琴}{lyre}；他也与他们在心神中一同舞蹈、一同喜乐，如同从前在天主的约柜前一般，\BibleRef{撒下 6:14} 他调和音乐和谐，以嗫嚅的声音甜蜜地咿呀说出：\Quote{奉主名而来的，当受赞颂。} 我们要问谁呢？先知，请告诉我们，这位奉主名而来的，是谁？祂会说：今天不是我教导你们的时候，因为祂已将学校奉献给幼童，祂从婴孩和吃奶者的口中完备了赞美，为毁灭仇敌和复仇者，好藉这些孩童的奇迹，使父亲的心转向儿女，使悖逆者转向义人的智慧。\BibleRef{拉 4:6} 那么，孩子们，请告诉我们，你们这优美而雅致的歌唱比赛从何而来？谁教给你们的？谁指导的？谁把你们聚集在一起？你们的书板是什么？你们的老师是谁？他们说：\Quote{你们只要加入我们，成为这歌声和庆节的同伴，就会知道那些梅瑟和先知切望之事。}\BibleRef{路 10:24} 既然孩子们邀请了我们，并与我们行了右手相交之礼，\BibleRef{迦 2:9} 让我们来吧，蒙爱者，我们自己也要效法那神圣的歌咏团，与宗徒们一起，为那位升上诸天之上、向着东方，且甘愿在地上骑着一头驴驹的主预备道路。让我们与孩子们一同高扬树枝，以橄榄枝欢喜鼓掌，好使圣神也吹在我们身上，使我们有序地高唱那由天主教导的曲调：\Quote{奉主名而来的，当受赞颂；贺三纳于至高之天。}\BibleRef{玛 21:5} 今天，圣祖雅各伯也在心神中庆祝，见到他的预言得以应验，并与信友一同朝拜圣父，看见那位将驴驹拴在葡萄树\BibleRef{创 49:10}上的主，骑着一头驴驹。今天，驴驹已预备好，这是外邦人——先前无理性者——的无理性象征，为表明外邦人的顺服；婴孩则宣告他们先前在认识天主方面的孩童状态，以及后来因敬拜天主与实践真宗教而达致的成全。今天，依照先知的话，光荣的君王在地上受光荣，并使我们地上的居民参加天上的筵席，为显示祂是二者的主，正如祂受到二者共同的赞美。因此，天上的万军歌唱，宣布地上的救恩：\Quote{圣、圣、圣，万军的上主！祂的光荣充满大地。}\BibleRef{依 6:3} 而地上的人，与天上欢乐的赞美诗和谐相应，高呼：\Quote{贺三纳于至高之天；贺三纳于达味之子！} 在天上，\OriginalTerm{赞颂词}{doxology}被高举：\Quote{愿上主的光荣从祂的处所受赞颂；}\BibleRef{则 3:22} 而在地上，这赞颂被以这些话语应和：\Quote{奉主名而来的，当受赞颂。}\par

三、当这些事正在发生时，门徒们因所见的一切异能而欢欣，大声赞美天主说：\Quote{因上主之名而来的君王，应受赞颂！在天上有和平，在至高之处有光荣。}\BibleRef{路 19:37} 城市开始询问说：\Quote{这人是谁？}\BibleRef{玛 21:10} 激发了他们硬化而根深蒂固的嫉妒，反对主的荣耀。但当你听我说城市时，要明白是指会堂中古老而无序的群众。他们忘恩负义、恶毒地问：\Quote{这人是谁？}好像从未见过他们的恩主，以及那位因神圣奇迹——超越人力——而闻名显赫者；因为黑暗没有领会\BibleRef{若 1:5}那照耀在它上面的不落之光。因此，先知依撒意亚非常贴切地向他们呼喊说：\Quote{聋子啊，听吧！瞎子啊，看吧，好让你们看见！谁比我的仆人更瞎？谁比那管辖他们的人更聋？}\BibleRef{依 42:18}\Quote{你们屡次看见，却不留意；耳朵开了，却不听见。}\BibleRef{依 42:20} 看，亲爱的，这些话是多么准确！圣神自己预见未来，通过祂的圣人们预言未来的事，如同现在一样。因为这些忘恩负义的人看见了，并通过祂的奇迹触摸了行奇事的天主，却仍然不信。\EditorialNote{若 9} 他们看见一个生来瞎眼的人，向他们宣告那位恢复他视力的天主。他们看见一个瘫痪者，似乎与他的病弱成为一体，在主的一声令下从疾病中解脱。\BibleRef{若 5:5} 他们看见拉匝禄\Person{拉匝禄}，从死亡区域被放逐。\BibleRef{若 11:44} 他们听说祂在海面上行走。\BibleRef{玛 14:26} 他们听说那未经栽培而供应的酒；\BibleRef{若 2:7} 那在自发宴席上食用的饼；\BibleRef{若 6:11} 他们听说魔鬼被驱散，病人恢复健康。他们的大街小巷宣告祂的奇事；他们的道路向旅行者宣告祂的医治能力。全犹太\Place{犹太}充满了祂的恩惠；然而现在，当他们听到神圣的赞美时，他们询问：\Quote{这人是谁？}噢，这些假教师的疯狂！噢，不信的父亲！噢，愚蠢的长老！噢，无耻客纳罕\Place{客纳罕}的种子，而不是虔敬犹大\Place{犹大}的种子！\BibleRef{达 3:56} 孩子们承认他们的创造主，但他们的不信的父母说：\Quote{这人是谁？}那年轻而没有经验的世代赞美天主，而那些在邪恶中老去的人询问：\Quote{这人是谁？}吃奶的婴孩赞美祂的神性，而长者们口出亵渎；孩子们虔诚地献上赞美的祭献，而亵渎的司铎们不虔敬地愤怒。\BibleRef{玛 21:15}\par

四、噢，你们这些对义人的智慧悖逆不服从的人，\BibleRef{路 1:17} 将你们的心转向你们的子女。学习天主的奥秘：所发生的事本身就证明，那正是天主藉未经教导的舌头所歌颂的。你们查考圣经，\BibleQuote{你们查考圣经}正如你们从主所听见的，\BibleRef{若 5:39} 因为正是这些圣经为我作证，不要对这奇迹无知。听吧，没有恩宠的、忘恩负义的人们，先知匝加利亚\Person{匝加利亚}给你们带来了什么好消息。他说：\Quote{熙雍\Place{熙雍}女子，你应尽量喜乐！看，你的君王到你这里来；正义的、胜利的、谦逊的，骑在驴驹上。}\BibleRef{匝 9:9} 为什么你们拒绝喜乐？为什么当太阳照耀时，你们喜爱黑暗？为什么你们针对不可战胜的和平图谋战争？因此，如果你们是熙雍\Place{熙雍}的子孙，就和你们的子女一同欢庆。让你们子女的宗教服务成为你们喜乐的借口。从他们那里学习谁是他们的老师；谁召集了他们；教义从何而来；这新神学和旧预言是什么意思。如果没有人教导他们这些，而是他们自动发出赞美的诗歌，那么你们该认出天主的工程，正如法律上所记载的：\Quote{从婴孩和哺乳者的口中，你备妥了赞美。}因此，你们要加倍喜乐，因为你们成了这些孩子的父母，他们在天主的教导下，赞美他们长辈所不知的事。将你们的心转向你们的子女，\BibleRef{路 1:17} 不要闭眼抗拒真理。但如果你们依然如故，听了却不听见，看了却不看见，\BibleRef{依 6:10} 且无缘无故地与你们的子女唱反调，那么按照救主的话，他们将是审判你们的人。\BibleRef{玛 12:27} 因此，先知依撒意亚\Person{依撒意亚}也曾预言你们这事，说：\Quote{雅各伯\Person{雅各伯}现在不再羞愧，他的面容不再苍白。当他们看见我的孩子们做我的工作时，他们必因我而圣化我的名，并圣化雅各伯的圣者，且敬畏以色列\Place{以色列}的天主。那些心神迷误的，将获得认识；那些抱怨的，将学习服从；那些口齿不清的舌头，将学习讲和平。}你们看，愚昧的犹太人，先知从一开始讲论时，就因你们的不信向你们宣告了混乱。你们甚至从他那里学习，他如何宣告你们的子女所发出的、由天主启示的赞美诗歌，正如达味圣王\Person{达味}所预言的：\Quote{从婴孩和哺乳者的口中，你备妥了赞美。}因此，要么——如同应该的——你们认领你们子女的虔敬为己有，要么虔敬地将你们的子女交给我们。我们要和他们一同起舞，并为新的光荣齐声歌唱这由天主启示的诗歌。\par

五. 确实，年迈的西默盎曾遇见救主\BibleRef{路 2:29}，并抱在怀中，那创造世界的婴孩，且宣告祂是主与天主。但现在，取代愚昧的长老，孩童们如同西默盎那样迎接救主，不是用他们的手臂，而是在祂脚下铺展树枝，祝福那骑在驴驹上、如同坐在革鲁宾上的主天主：\Quote{贺三纳于达味之子！因主名而来的，当受赞美！}让我们也与他们一同高呼：\Quote{因主名而来的，当受赞美！}祂是光荣的君王天主，为了我们而变成贫穷，却仍拥有自己的尊荣，不知贫穷为何物，好藉祂的慷慨使我们富足。当受赞美者，祂曾谦卑而来，今后将要在荣耀中再来：第一次，卑微地骑在驴驹上，由婴孩们颂扬，为应验经上所记：\Quote{天主啊，祢的行径已被看见，我天主、我君王的行径，在圣所中。}但第二次，祂要坐在云上，威荣可畏，由天使与万有随从。啊，孩童们甘甜的舌音！啊，中悦天主者的真诚教诲！达味在预言中把神性隐藏在字句下；孩童们开启宝库，口舌中涌出财富，以满含恩宠的话语，清楚地邀请所有人享用。因此让我们与他们一同汲取那永不凋谢的财富。在我们永不满足的胸怀和无法填满的宝库中，收藏神圣的恩赐。让我们不停地高呼：\Quote{因主名而来的，当受赞美！}真天主，以真天主之名，全能者自全能者，圣子因圣父之名而来。真君王自真君王而来，祂的国度如同生父的国度，与永恒同在，与永恒同存且先于永恒。因为这是两者共有的；圣经并未将这份尊荣归于圣子，仿佛来自他处，或有开始，或可增减——绝无此意！——而是祂由本性、由真实而合宜的所有权所应得的。因为圣父、圣子及圣神的国度是唯一的，正如他们的\OriginalTerm{本体}{substance}是唯一的，统治也是唯一的。因此，我们以同一敬拜，敬拜\OriginalTerm{三位}{three Persons}之中的唯一神性，无始、无受造、无终、且无继承者。因为圣父绝不会不再是圣父，圣子也绝不会不再是圣子与君王，圣神也绝不会不再是其本体与位格之所是。因为\OriginalTerm{天主圣三}{Trinity}无论在永恒、共融或主权上，都不会有任何减损。圣子之所以被称为君王，并非因为祂为了我们降生成人，在肉身上击败了那反对我们的暴君，藉此获得了战胜其残忍仇敌的胜利，而是因为祂永远是主与天主；因此现在，既在祂取得肉体之后，也永远，祂仍是君王，如同生祂的那一位。异端者啊，不要反对基督的王国，以免羞辱了生祂的那一位。你若是有信德的，就当以信德亲近基督，我们的真天主，而不是把自由当作邪恶的托辞。你若是仆人，就当战战兢兢地服从你的主人；因为那与圣言作战的，不是好仆人，而是明显的仇敌，正如经上所记：\Quote{不尊敬子的，就是不尊敬派遣祂来的父。}\par

六. 但是，亲爱的，让我们在讨论中回到我们离题之处，高呼：\Quote{因主名而来的，当受赞美！}那位良善仁慈的牧人，甘愿为羊舍命。正如猎人用羊捕获吞吃羊的狼，同样，大司牧\BibleRef{伯前 5:4}将自己作为人献给那属灵的、毁灭灵魂的狼，好藉那个曾被它们掠食的亚当，使掠食者成为祂的掠物。因主名而来的，当受赞美：天主对抗魔鬼；并非显明在祂不可逼视的威能中，而是借肉身的软弱，捆绑那对敌我们的壮士\BibleRef{玛 12:29}。因主名而来的，当受赞美：君王对抗暴君；不是以全能的力量与智慧，而是以被视为愚拙的十字架\BibleRef{格前 1:21}，这十字架从那在邪恶中狡猾的蛇身上夺去了它的掠物。因主名而来的，当受赞美：真理对抗说谎者；救主对抗毁灭者；和平的君王\BibleRef{依 9:6}对抗挑起战争者；爱人的天主对抗恨人者。因主名而来的，当受赞美：主怜悯祂手中的受造物。因主名而来的，当受赞美：主拯救那迷途的人；驱除错谬；光照黑暗中的人；废除偶像的欺骗；代之以对天主的救恩认识；圣化世界；驱逐假神崇拜的可憎与悲惨。因主名而来的，当受赞美：那一位为众人；拯救穷人脱离强于他者的手，甚至穷苦贫乏者脱离那劫掠他的。因主名而来的，当受赞美：向落入强盗手中\BibleRef{路 10:34}且被弃之不顾的人倒上酒和油。因主名而来的，当受赞美：拯救我们，不是藉着使者或天使，而是主亲自拯救了我们\BibleRef{依 63:9}。因此我们也赞美你，主；你与圣父及圣神，在万世之前及永远，当受赞美。确实，在万世之前及直到如今，你本是无形无体的，但从现在直到永远，你拥有那不能改变的神圣人性，且你永不与它分离。\par

七. 让我们再看看随后发生的事。最神圣的福音作者怎么说？当主进入圣殿时，盲人和跛者来到祂跟前；祂治好了他们。当司祭长和法利塞人看见祂所做的奇事，并听见儿童呼喊说：\BibleQuote{贺三纳归于达味之子！奉主名而来的当受赞美！}\BibleRef{玛 21:14}他们不能容忍人们对祂的敬意，就来到祂跟前这样说：\Quote{你没有听见他们说什么吗？}好像他们说：\Quote{你听见这些无辜者说出只属于天主、唯独属于天主的事，难道不气愤吗？天主岂不早就藉先知晓谕：\BibleQuote{我的荣耀我必不给予别人；}\BibleRef{依 42:8}而你作为人，怎么使自己成为天主呢？}\BibleRef{若 10:33}但是长期忍耐的那一位，祂是慈悲丰厚的，\BibleRef{岳 2:13}且缓于发怒的，\BibleRef{雅 1:18}如何回答此事？祂容忍这些狂乱的人；用辩词抑制他们的怒气；祂转而提醒他们圣经；祂为所发生的事提出见证，并不回避质询。因此祂说：\BibleQuote{你们从未听过我藉先知所说：那时你们就知道我是那个发言的？}\BibleRef{依 52:6}又：\Quote{从婴孩和吃奶者的口中，你因仇敌的缘故成全了赞美，为使你摧毁仇敌和报复者？}无疑你们就是这些人，你们遵循法律，阅读先知，却轻视我，而我在法律和先知中早被预言了。你们确实以虔诚为借口，想要报复天主的荣耀，却不明白那鄙视我的也鄙视我的父。\BibleRef{若 15:23}\BibleQuote{我从天主而出，来到世上，}\BibleRef{若 16:28}\BibleQuote{我的荣耀也是我父的荣耀。}就这样，这些愚妄的人，被我们的救主天主说服，真理堵住了他们的嘴，便不再回答祂；却采取了新的愚妄计谋，商议反对祂。但让我们歌唱：\Quote{我们的主何其伟大，祂的能力何其浩大；祂的智慧不可胜数。}所有这一切是为了让天主的羔羊和圣子，祂除免世罪，自愿为我们走向祂救赎的苦难，并好像在市场和售卖处被认出；为了让那些买祂的人用三十块银钱为祂立约，而祂要以赋予生命的血救赎世界；为了让基督，我们的逾越节，为我们而被祭杀，以便那些被祂宝血所洒、嘴唇被盖印如同门框的\BibleRef{出 11:7}能逃脱毁灭者的箭矢；并为了让基督在肉身中这样受苦，第三日复活，与圣父和圣神享有同等的尊威和荣耀，被一切受造物同等敬拜；因为\BibleQuote{天上、地上和地下的万膝都将向祂跪拜，}\BibleRef{斐 2:10}把荣耀归于祂，至于无穷之世。阿们。\par

\SourceRecord{Translated by William R. Clark.；Ante-Nicene Fathers；Vol. 6.；Edited by Alexander Roberts, James Donaldson, and A. Cleveland Coxe.；Buffalo, NY: Christian Literature Publishing Co.,；1886.；<http://www.newadvent.org/fathers/0628.htm>.}
